%! Polar Function Graphs — Unit 7, Lesson 7.7 — Homework
\documentclass[10pt]{article}
\usepackage{precalculus-article}
\usepackage{precalculus-boxes}
\newcommand{\TallMath}[1]{$\displaystyle #1\rule[-1.4em]{0pt}{3.2em}$}

\begin{document}

\pageheader{Unit 7, Lesson 7.7}{Homework}
\namedateperiod

\vspace{0.10in}

\begin{notesbox}{Problems 1--6}

\textbf{1.}\quad The table below gives values of $r = 1 - \sin\theta$ for selected
angles. Use it to answer the questions.

\vspace{4pt}
\renewcommand{\arraystretch}{1.55}
\begin{tabular}{|c|c|c|c|c|c|c|c|}
\hline
$\theta$ & $0$ & $\dfrac{\pi}{6}$ & $\dfrac{\pi}{2}$ & $\pi$ &
           $\dfrac{3\pi}{2}$ & $\dfrac{5\pi}{4}$ & $2\pi$ \\
\hline
$r = 1 - \sin\theta$ & $1$ & $0.5$ & $0$ & $1$ & $2$ & $1.71$ & $1$ \\
\hline
\end{tabular}

\vspace{6pt}
(a)\quad What is the maximum value of $r$? \underline{\hspace{1.5cm}}
At what angle does this occur? \underline{\hspace{2cm}}

\vspace{4pt}
(b)\quad At what angle is $r = 0$? \underline{\hspace{2cm}}
The curve passes through the \underline{\hspace{2cm}} at that angle.

\vspace{4pt}
(c)\quad On the interval $\theta \in \left[0, \dfrac{\pi}{2}\right]$, is $r$
increasing or decreasing? \underline{\hspace{2.5cm}}

\vspace{0.14in}
\textbf{2.}\quad The rectangular graph of $r = \cos(3\theta)$ is shown below for
$\theta \in [0, \pi]$. Use it to determine how many petals the polar curve has.

\vspace{4pt}
\begin{center}
\begin{tikzpicture}[x=1.55cm, y=1.1cm]
  \draw[linegray, very thin] (-0.2,-1.3) grid[xstep=0.25, ystep=0.5] (3.3,1.3);
  \draw[->] (-0.2,0) -- (3.45,0) node[right]{\small $\theta$};
  \draw[->] (0,-1.35) -- (0,1.4) node[above]{\small $r$};
  \foreach \x/\lbl in {0.524/$\frac{\pi}{6}$, 1.047/$\frac{\pi}{3}$,
                        1.571/$\frac{\pi}{2}$, 2.094/$\frac{2\pi}{3}$,
                        2.618/$\frac{5\pi}{6}$, 3.142/$\pi$} {
    \draw (\x,0.07) -- (\x,-0.07);
    \node[below, font=\scriptsize] at (\x,-0.07) {\lbl};
  }
  \foreach \y/\lbl in {-1/{$-1$}, 1/{$1$}} {
    \draw (0.07,\y) -- (-0.07,\y);
    \node[left, font=\scriptsize] at (-0.07,\y) {\lbl};
  }
  \draw[navy, thick, domain=0:3.14, samples=200]
    plot (\x, {cos(3*deg(\x))});
  \node[right, font=\small, navy] at (3.2, 0.5) {$r = \cos(3\theta)$};
\end{tikzpicture}
\end{center}

\vspace{4pt}
(a)\quad List all $\theta$ values in $[0, \pi]$ where $r = 0$:
\underline{\hspace{5cm}}

\vspace{4pt}
(b)\quad How many complete zero-to-max/min-to-zero arcs are visible in $[0, \pi]$?
\underline{\hspace{1.5cm}}

\vspace{4pt}
(c)\quad The formula says $r = \cos(n\theta)$ with $n$ odd produces $n$ petals.
For $n = 3$, the polar curve $r = \cos(3\theta)$ has \underline{\hspace{1cm}} petals.

\vspace{0.14in}
\textbf{3.}\quad (Domain Restriction) Consider $r = \sin(2\theta)$.

\vspace{4pt}
(a)\quad When $\theta \in \left[0, \dfrac{\pi}{2}\right]$, how many petals of
$r = \sin(2\theta)$ are traced? \underline{\hspace{1.5cm}}

\vspace{4pt}
(b)\quad Describe the portion of the curve traced when $\theta \in [0, \pi]$:
\writelines{2}

\vspace{0.14in}
\textbf{4.}\quad (AP-Style MC) For the polar function $r = 2 + 2\sin\theta$, which
of the following is the maximum distance from the origin?

\begin{enumerate}[label=(\Alph*), leftmargin=2em, itemsep=3pt]
  \item $2$
  \item $4$
  \item $2\pi$
  \item $1$
\end{enumerate}

Answer: \underline{\hspace{1.5cm}}

\vspace{0.14in}
\textbf{5.}\quad (AP-Style MC) The polar curve $r = \sin(4\theta)$ is a rose curve
with $n = 4$. Since $n$ is even, the number of petals is $2n$. How many petals does
this curve have?

\begin{enumerate}[label=(\Alph*), leftmargin=2em, itemsep=3pt]
  \item $2$
  \item $4$
  \item $8$
  \item $16$
\end{enumerate}

Answer: \underline{\hspace{1.5cm}}

\vspace{0.14in}
\textbf{6.}\quad The function $r = 3$ is a polar function with a constant output.

\vspace{4pt}
(a)\quad For every angle $\theta$, the radius is always \underline{\hspace{1.5cm}}.

\vspace{4pt}
(b)\quad What shape does the polar curve $r = 3$ trace? \underline{\hspace{4cm}}

\vspace{4pt}
(c)\quad What is the domain of $\theta$ needed to trace the complete curve once?
\underline{\hspace{4cm}}

\end{notesbox}

\vspace{0.08in}

\begin{tcolorbox}[colback=navylight, colframe=navy, arc=2mm,
    left=3mm, right=3mm, top=2mm, bottom=2mm]
\small\textbf{Preview --- Lesson 3.15: Rates of Change of Polar Functions}\\
You can now read a polar function from tables and rectangular graphs. In Lesson 3.15
we'll ask: how fast is the radius $r$ changing as the angle $\theta$ increases?
Come ready to connect the slope of the rectangular graph to the behavior of the
polar curve.
\end{tcolorbox}

\end{document}
